\section{Datos}\label{sec:datos}

Se usan tres conjuntos públicos de Hugging Face, uno por tipo de tarea (tabla~\ref{tab:datos}).

\begin{table}[ht]
\centering
\caption{Conjuntos de datos y uso en el proyecto.}\label{tab:datos}
\small
\begin{tabularx}{\linewidth}{@{}l l X c@{}}
\toprule
Dataset & Partición & Uso & Consultas \\
\midrule
WikiQA~\cite{wikiqa} & \code{validation} & Base de conocimiento en ChromaDB y preguntas abiertas & 20 \\
GSM8K~\cite{gsm8k} & \code{test} & Problemas matemáticos de varios pasos; referencia numérica tras \code{\#\#\#\#} & 15 \\
MBPP~\cite{mbpp} & \code{test} & Generación de funciones Python; verificación con \code{test\_list} & 15 \\
\midrule
\multicolumn{3}{@{}l}{Total} & 50 \\
\bottomrule
\end{tabularx}
\end{table}

\paragraph{Base de conocimiento.}
WikiQA asocia cada pregunta con varias frases candidatas de Wikipedia, marcadas como correctas (\code{label=1}) o incorrectas (\code{label=0}).
Se agrupan las filas por pregunta y se conservan las 126 preguntas de validación con al menos una frase correcta.
Se indexan \emph{todas} sus frases candidatas: 1119 frases únicas, también las incorrectas.
Indexar solo las correctas pondría la respuesta al alcance de cualquier recuperación y favorecería artificialmente al Naive RAG.
Con las incorrectas, la base contiene distractores temáticamente cercanos.
El índice es una colección persistente de ChromaDB con su función de embeddings por defecto.
Los identificadores salen de un hash del contenido, así que reindexar no duplica documentos.

\paragraph{Consultas de MBPP.}
Siguiendo a Austin et al.~\cite{mbpp}, la pregunta incluye los \code{assert} de \code{test\_list} para que el modelo conozca el nombre y la firma de la función.
Ambos métodos ven esas pruebas, pero solo el multi-agente puede ejecutarlas.
Su exactitud en MBPP mide, por lo tanto, si sabe verificar y corregir su código, no si generaliza a pruebas ocultas.
